\documentclass[a4paper,11pt]{article}
\usepackage{exam-style}

\fancyhead[L]{\fontsize{9pt}{10pt}\selectfont\color{ctp-subtext0}341 农业综合知识三（信电）· 模拟卷七}
\fancyhead[R]{\fontsize{9pt}{10pt}\selectfont\color{ctp-subtext0}信电学院部分}

\begin{document}
\examtitle

% ============================================================
% 第一部分：程序设计基础知识（75 分）
% ============================================================
\parttitle{第一部分 \quad 程序设计基础知识（共75分）}

% ---------- 一、填空题（10分） ----------
\qtypename{一、填空题}{共5题，每题2分，共10分}
\anslines{0.1cm}

\q 已知 \code{int a = 0, b = 1, c = 2;}，则表达式 \code{a++ \&\& b++ || c++} 求值后，\code{a}、\code{b}、\code{c} 的值分别是 \fillblank 、 \fillblank 、 \fillblank 。

\q 若有 \code{char s[] = "012\textbackslash{}034\textbackslash{}0ABC";}，则 \code{strlen(s)} 的值是 \fillblank ，\code{sizeof(s)} 的值是 \fillblank 。

\q 已知 \code{int a[3][2] = \{\{1,2\},\{3,4\},\{5,6\}\};}，则 \code{*(a[0] + 1)} 的值是 \fillblank ，\code{*(*(a + 2) + 1)} 的值是 \fillblank 。

\q 若有 \code{typedef int (*ArrPtr)[5];} 则 \code{ArrPtr p;} 定义了一个 \fillblank 。

\q 执行 \code{fseek(fp, 0L, SEEK\_SET);} 后，文件指针指向 \fillblank 。

% ---------- 二、简答题（15分） ----------
\qtypename{二、简答题}{共5题，每题3分，共15分}

\q 请解释 C 语言中 \code{const int *p}、\code{int *const p} 和 \code{const int *const p} 三者的区别，并举例说明各自的使用场景。
\anslines{1.5cm}

\q 什么是"内存泄露"？写出一个会导致内存泄露的代码示例，并说明如何避免。
\anslines{1.5cm}

\q 解释函数指针的用途，并写出一个使用函数指针实现回调的简单示例。
\anslines{1.5cm}

\q 说明 \code{\#include <file>} 与 \code{\#include "file"} 在搜索路径上的区别。当头文件名与标准库头文件同名时会发生什么？
\anslines{1.5cm}

\q 简述 \code{typedef} 与 \code{\#define} 在定义类型别名时的本质区别，并举例说明 \code{\#define} 可能带来什么问题。
\anslines{1.5cm}

% ---------- 三、分析题（20分） ----------
\qtypename{三、分析题}{共10题，每题2分，共20分}

阅读下列程序片段，写出运行结果。

\q
\begin{lstlisting}[language={[custom]C}]
int a = 0, b = 1;
if (a++ && b++)
    printf("T");
else
    printf("F");
printf(" %d %d", a, b);
\end{lstlisting}
运行结果：\underline{\hspace{3cm}}


\q
\begin{lstlisting}[language={[custom]C}]
int x = 0;
while (x++ < 5);
printf("%d", x);
\end{lstlisting}
运行结果：\underline{\hspace{3cm}}


\q
\begin{lstlisting}[language={[custom]C}]
int a[] = {1, 2, 3, 4, 5};
int *p = a + 2;
printf("%d ", p[-1]);
printf("%d", *(a + 4));
\end{lstlisting}
运行结果：\underline{\hspace{3cm}}


\q
\begin{lstlisting}[language={[custom]C}]
int i, j;
for (i = 0; i < 3; i++)
    for (j = 0; j < 3; j++)
        if (i == j)
            continue;
        else
            goto end;
end:
printf("i=%d j=%d", i, j);
\end{lstlisting}
运行结果：\underline{\hspace{3cm}}


\q
\begin{lstlisting}[language={[custom]C}]
int a = 1, b = 2;
a ^= b; b ^= a; a ^= b;
printf("%d %d", a, b);
\end{lstlisting}
运行结果：\underline{\hspace{3cm}}


\q
\begin{lstlisting}[language={[custom]C}]
int func(int x) {
    static int n = 0;
    return x + (n++);
}
int main() {
    printf("%d ", func(5));
    printf("%d", func(5));
    return 0;
}
\end{lstlisting}
运行结果：\underline{\hspace{3cm}}


\q
\begin{lstlisting}[language={[custom]C}]
char *s = "Hello";
printf("%c ", *(s + 1));
printf("%c", s[4]);
\end{lstlisting}
运行结果：\underline{\hspace{3cm}}


\q
\begin{lstlisting}[language={[custom]C}]
int x = 10;
{
    int x = 20;
    printf("%d ", x);
}
printf("%d", x);
\end{lstlisting}
运行结果：\underline{\hspace{3cm}}


\q
\begin{lstlisting}[language={[custom]C}]
int a[5] = {0};
for (int i = 0; i < 5; i++)
    a[i] = i * 2 - 1;
printf("%d", a[2] + a[4]);
\end{lstlisting}
运行结果：\underline{\hspace{3cm}}


\q
\begin{lstlisting}[language={[custom]C}]
#include <string.h>
char s[20] = "ABC";
strcat(s, "DEF");
printf("%zu %zu", strlen(s), sizeof(s));
\end{lstlisting}
运行结果：\underline{\hspace{3cm}}

% ---------- 四、程序设计题（30分） ----------
\qtypename{四、程序设计题}{共2题，每题15分，共30分}

\pq 编写 C 语言程序：定义函数 \code{void matrix_multiply(int A[][3], int B[][2], int C[][2], int m, int n, int p)}，实现 $m \times n$ 矩阵 A 与 $n \times p$ 矩阵 B 的乘法，结果存入 $m \times p$ 矩阵 C。主函数中定义两个矩阵（\code{A[2][3]}, \code{B[3][2]}），从键盘输入元素值，调用函数计算乘积并输出结果矩阵。要求处理输入验证（如维度不匹配时给出提示）。

\anslines{5cm}

\pq 编写 C 语言程序：定义单向链表结构 \code{struct ListNode \{int val; struct ListNode *next;\};}，实现以下函数：
\begin{itemize}
  \item \code{struct ListNode* createList(int arr[], int n);} —— 根据数组创建链表；
  \item \code{struct ListNode* reverseList(struct ListNode *head);} —— 原地反转链表；
  \item \code{void printList(struct ListNode *head);} —— 打印链表。
\end{itemize}
在主函数中从键盘输入一组整数（以 -1 结束），创建链表，反转后输出。要求处理空链表情况。

\anslines{6cm}

% ============================================================
% 第二部分：数据库技术与应用（75 分）
% ============================================================
\parttitle{{第二部分 \quad 数据库技术与应用（共75分）}}

% ---------- 一、填空题（10分） ----------
\qtypename{一、填空题}{共5题，每题2分，共10分}

\q 设有关系 R(A, B, C) 和 S(B, C, D)，则自然连接 R $\bowtie$ S 的结果属性集为 \fillblank ，元组需满足的条件是 \fillblank 。

\q 关系代数中，完成"从表中选取满足条件的元组"的操作称为 \fillblank 运算，完成"从表中选取若干属性列"的操作称为 \fillblank 运算。

\q 在 MySQL 中，\code{UNION} 与 \code{UNION ALL} 的区别在于 \fillblank 。

\q 数据库并发操作中，"脏读"是指 \fillblank ，"幻读"是指 \fillblank 。

\q 在 MySQL 中，\code{TRUNCATE TABLE} 与 \code{DELETE FROM} 的区别是：\fillblank 是 DDL 语句，\fillblank 是 DML 语句；前者 \fillblank 回滚，后者 \fillblank 回滚。

% ---------- 二、简答题（10分） ----------
\qtypename{二、简答题}{共2题，每题5分，共10分}

\q 什么是数据依赖？函数依赖与多值依赖有何区别？请分别举例说明。在关系规范化中，多值依赖导致的冗余问题需要通过什么范式解决？
\anslines{3cm}

\q 简述数据库并发控制中"两段锁协议（2PL）"的基本规则。2PL 如何保证事务的可串行化调度？什么是"严格两段锁协议"？它对级联回滚问题有何改善？
\anslines{3cm}

% ---------- 三、操作题（15分） ----------
\qtypename{三、操作题}{本题共15分}

设有图书借阅管理系统关系模式：\\
读者 R(\underline{Rid}, Rname, Rtype, Dept, MaxBorrow)\\
图书 B(\underline{Bid}, Bname, Author, Publisher, Price, Total, Stock)\\
借阅 L(\underline{Rid, Bid, BorrowDate}, DueDate, ReturnDate, RenewCount)

说明：\code{Total} 为馆藏总量，\code{Stock} 为当前可借数量，\code{RenewCount} 为续借次数。

\q （2分）查询当前借阅未还（ReturnDate IS NULL）且已超过应还日期（DueDate）的读者编号、姓名及超期天数。
\anslines{1cm}

\q （2分）查询借阅次数最多的前 10 本书的书名、作者及借阅总次数。
\anslines{1cm}

\q （2分）查询每位读者当前借阅图书的数量，输出读者编号、姓名和借阅数量，仅显示借阅数量超过 3 的读者。
\anslines{1cm}

\q （2分）查询没有被任何读者借阅过的图书编号和书名（用 NOT EXISTS）。
\anslines{1cm}

\q （2分）用 SQL 统计每个出版社的图书总量和平均价格，仅显示图书总量大于 5 的出版社。
\anslines{1cm}

\q （2分）查询与"张三"借阅过完全相同图书的读者编号和姓名（即借阅图书集合相同）。
\anslines{1.2cm}

\q （3分）创建触发器 \code{trg\_return\_book}，当在 L 表更新 \code{ReturnDate} 字段（即还书操作）时，自动将对应图书的 Stock 加 1。
\anslines{1.5cm}

% ---------- 四、分析题（10分） ----------
\qtypename{四、分析题}{共1题，共10分}

\q 设有关系模式 R(A, B, C, D, E)，函数依赖集 F = \{AB $\rightarrow$ C, C $\rightarrow$ D, D $\rightarrow$ B, D $\rightarrow$ E\}。

\begin{enumerate}[leftmargin=2em, itemsep=3pt]
  \item[（1）] （4分）求 R 的所有候选键，要求写出完整推导过程。
  \item[（2）] （3分）R 最高属于第几范式？说明理由。
  \item[（3）] （3分）将 R 无损分解到 BCNF。
\end{enumerate}

\anslines{3.5cm}

% ---------- 五、设计题（10分） ----------
\qtypename{五、设计题}{共1题，共10分}

\q 某在线教育平台需要设计一个课程学习管理系统，已知语义：
\begin{itemize}
  \item 每位学生有学号、姓名、性别、入学年份、专业；
  \item 每位教师有工号、姓名、职称、所属学院；
  \item 每门课程有课程编号、课程名称、学分、学时、课程性质（必修/选修）；
  \item 每门课程由一位教师主讲，一位教师可以主讲多门课程；
  \item 学生可选修多门课程，每门课程可被多名学生选修，选修时需记录成绩和评语；
  \item 课程可以设置先修课程关系（即一门课程可以有若干门先修课程）；
  \item 每门课程可以安排多次考试（期中、期末、补考等），每次考试有考试编号、考试名称、考试时间、考试地点和最高分、平均分统计。
\end{itemize}
请完成：
\begin{enumerate}[leftmargin=2em, itemsep=3pt]
  \item[（1）] （5分）画出该系统的 E-R 图。
  \item[（2）] （5分）将 E-R 图转换为关系模式（标明主键和外键）。
\end{enumerate}

\anslines{4.5cm}

% ---------- 六、应用题（20分） ----------
\qtypename{六、应用题}{共2题，每题10分，共20分}

\q 现有 MySQL 数据库中的银行账户系统表：
\begin{lstlisting}[language=SQL]
CREATE TABLE accounts (
    acc_id     INT PRIMARY KEY,
    cust_name  VARCHAR(50) NOT NULL,
    balance    DECIMAL(12,2) NOT NULL DEFAULT 0.00,
    status     VARCHAR(10) DEFAULT 'active'
);

CREATE TABLE transactions (
    tx_id      INT PRIMARY KEY AUTO_INCREMENT,
    acc_id     INT NOT NULL,
    tx_type    VARCHAR(10) NOT NULL,  -- 'deposit','withdraw','transfer'
    amount     DECIMAL(12,2) NOT NULL,
    tx_time    DATETIME DEFAULT CURRENT_TIMESTAMP,
    status     VARCHAR(10) DEFAULT 'pending',
    FOREIGN KEY (acc_id) REFERENCES accounts(acc_id)
);
\end{lstlisting}
\begin{enumerate}[leftmargin=2em, itemsep=3pt]
  \item[（1）] （3分）查询所有余额低于 0 的账户信息（理论上余额不应为负——说明可能存在数据异常），输出账户编号、姓名和余额。
  \item[（2）] （3分）编写 SQL 语句，查询每个账户的最近 5 笔交易记录（按时间降序），输出账户编号、交易类型、金额和交易时间。
  \item[（3）] （4分）编写存储过程 \code{transfer\_funds(IN from\_acc INT, IN to\_acc INT, IN amount DECIMAL(12,2))}，实现账户间转账：检查两个账户是否存在且状态为 active，检查转出账户余额是否充足，执行转账并记录两条交易记录（一条转出，一条转入）。要求使用事务确保原子性，并在出错时回滚且输出错误信息（如余额不足、账户不存在等）。
\end{enumerate}
\anslines{3.5cm}

\q 现有 MySQL 数据库，回答以下问题：
\begin{enumerate}[leftmargin=2em, itemsep=3pt]
  \item[（1）] （3分）现有表 \code{employee(eid, ename, salary, dno)}，为常用查询 \code{SELECT * FROM employee WHERE dno = ? ORDER BY salary DESC} 设计合适的索引，并说明为什么这种索引能提高查询效率。
  \item[（2）] （3分）什么是 MySQL 的"覆盖索引"（Covering Index）？覆盖索引为什么能减少查询的 I/O 开销？请举例说明。
  \item[（3）] （4分）请设计一个存储过程 \code{backup\_table(IN tbl\_name VARCHAR(64))}，使用游标遍历指定表的所有列信息（从 \code{INFORMATION\_SCHEMA.COLUMNS} 获取），动态生成 \code{CREATE TABLE ... LIKE} 语句来备份表结构。说明使用动态 SQL（\code{PREPARE} / \code{EXECUTE}）的必要性。
\end{enumerate}
\anslines{3.5cm}

\end{document}
